\par\SourceLeaf{391}%
Die zweiten Differentialquotienten von $q$ sollen $p^{(2)},p^{(3)},p^{(4)}$ sein. Diess vorausgesetzt, sind $A,B,C$ Functionen von $q$, und so setzen wir
\[
\frac{dA}{dq}=\mathfrak{A}',\qquad\frac{d\mathfrak{A}'}{dq}=\mathfrak{A}'',\qquad\frac{dB}{dq}=\mathfrak{B}',\quad\text{etc.}
\]
und haben
\[
\begin{gathered}
A'=p'\mathfrak{A}',\qquad A''=p''\mathfrak{A}',\qquad B'=p'\mathfrak{B}',\quad\text{etc.}\\
\begin{aligned}
A^{(2)}&=p'p'\mathfrak{A}''+p^{(2)}\mathfrak{A}'\\
A^{(3)}&=p'p''\mathfrak{A}''+p^{(3)}\mathfrak{A}'\\
A^{(4)}&=p''p''\mathfrak{A}''+p^{(4)}\mathfrak{A}'.
\end{aligned}
\end{gathered}
\]
[Dadurch gehen die Gleichungen 12 und 13 über in die folgenden]
\[
\begin{aligned}
na&=p'\mathfrak{A}'+p''(C\mathfrak{B}'-B\mathfrak{C}')\\
nb&=p'\mathfrak{B}'+p''(A\mathfrak{C}'-C\mathfrak{A}')\\
nc&=p'\mathfrak{C}'+p''(B\mathfrak{A}'-A\mathfrak{B}')\\[4pt]
n\alpha&=p''\mathfrak{A}'-p'(C\mathfrak{B}'-B\mathfrak{C}')\\
n\beta&=p''\mathfrak{B}'-p'(A\mathfrak{C}'-C\mathfrak{A}')\\
n\gamma&=p''\mathfrak{C}'-p'(B\mathfrak{A}'-A\mathfrak{B}').
\end{aligned}
\]
[Aus diesen Relationen folgt auch, da $m=1$ ist:]
\[
nn=p'p'+p''p''.
\]
\par\indent Die obigen beiden Formeln $\odot$ geben jetzt [da $m'=m''=0$ ist]:
\[
\begin{aligned}
0&=p''p^{(3)}-p'p^{(4)}+p''(p'p'+p''p'')\Sigma A(\mathfrak{B}'\mathfrak{C}''-\mathfrak{C}'\mathfrak{B}'')\\
0&=p''p^{(2)}-p'p^{(3)}+p'(p'p'+p''p'')\Sigma A(\mathfrak{B}'\mathfrak{C}''-\mathfrak{C}'\mathfrak{B}'')
\end{aligned}
\]
und hieraus
\[
p''p''p^{(2)}-2p''p'p^{(3)}+p'p'p^{(4)}=0
\]
und hieraus auch
\[
A''A''A^{(2)}-2A''A'A^{(3)}+A'A'A^{(4)}=0
\]
etc.
